% ============================================================
% qp-headfoot.tex —— 页眉页脚模块（全品式导学件，供 main.tex \input）
% 职责：页脚 twoside 奇偶（奇＝右端黑粗小字「章名＋件名」＋页码灰块；偶＝左端页码灰块＋黑粗小字册名）。
%       v4.2 撤页眉（翻拍板8②回归原旨：全品 p04-p06 实证无页眉；geometry headheight/headsep=0pt，
%       版心顶＝15mm 贴全品）；\fancyhead 调用撤、\qpjieming 宏保留作他件复用参数。
% 参数宏（他件复用改这四条即可）：
%   \qpzhangming 章名　\qpjieming 页眉节名（v4.2 起不排印，保留备他件）　\qpjianming 件名　\qpceming 册名
% v4.2 页脚小字：全品页脚小字实证灰度 58（近黑）且加粗——\qphuitext 由 6.5pt midgray 改 6.5pt 黑加粗；
% 页码块数字黑粗维持（口径不变）。页码块几何：24mm×7mm，底缘离页底 11mm（拍板22 保持）；数字水平+垂直居中——
%   块相对页脚基线 −2.275..+4.725mm（strut 负深度），数字 cap 中心 +1.225mm＝块中心；
%   \raisebox [0pt][0pt] 防 fancyhdr 盒高重算（同前轮 A6 方案）。
% v4.2 拍板22 让位改排：bottom=20mm 后 footskip=13pt 自然页脚基线＝版心底下 13pt＝15.41mm（离页底），
%   自然块底 15.41−2.275＝13.13mm——压不到 11mm 目标，基线整体降 2.11mm（raisebox −2.11mm，块底 11.02mm）。
% 调用关系：main.tex → qp-layout（geometry/颜色）→ 本模块 → qp-titles → qp-blocks。
% ============================================================
\usepackage{fancyhdr}

% ---------- 参数宏（他件复用入口） ----------
\newcommand{\qpzhangming}{第一章\quad 空间向量与立体几何}
\newcommand{\qpjieming}{1.1.1 空间向量及其运算}   % v4.2 撤页眉后不排印，保留作他件复用参数
\newcommand{\qpjianming}{导学件}
\newcommand{\qpceming}{高中数学\quad 选择性必修第一册（人教B版）}

% ---------- 页码灰块（底 DDDDDD＋黑加粗数字，居中） ----------
\newcommand{\qpshuzi}{\setlength{\fboxsep}{0pt}\colorbox{pnumbg}{%
  \makebox[24mm][c]{\rule[-2.275mm]{0pt}{7mm}\fontsize{10.5pt}{12pt}\selectfont\bfseries\thepage}}}
% v4.2：页脚小字＝6.5pt 黑加粗（全品实测灰度 58 近黑且粗；原 midgray 灰撤）
\newcommand{\qphuitext}{\fontsize{6.5pt}{8pt}\selectfont\bfseries}

\pagestyle{fancy}
\fancyhf{}
\renewcommand{\headrulewidth}{0pt}
% v4.2 撤页眉：\fancyhead[C]{\qphuitext\qpjieming} 整行撤（翻拍板8②回归原旨），页眉带零内容零占位
\fancyfoot[RO]{\raisebox{-2.11mm}[0pt][0pt]{\qphuitext\qpzhangming\quad\qpjianming\hspace{3mm}\qpshuzi}}
\fancyfoot[LE]{\raisebox{-2.11mm}[0pt][0pt]{\qpshuzi\hspace{3mm}\qphuitext\qpceming}}
